\documentclass[12pt]{article}
\usepackage[T2A]{fontenc}
\usepackage[utf8]{inputenc}
\usepackage[russian]{babel}
\usepackage{latexsym,amssymb,amsthm}
\usepackage{amsmath}
\usepackage{wrapfig}
\RequirePackage[pdftex]{graphicx}
%%\usepackage{array}

\usepackage{epsf,epic,eepic}
\input ris.tex
% \input cells.sty
% \addrisoption{\risleftfalse}

\DeclareRobustCommand{\No}{\ifmmode{\nfss@text{\textnumero}}\else\textnumero\fi}

\oddsidemargin=0pt
\textwidth=159mm
\headheight=0pt
\headsep=0pt
\topmargin=0pt
\textheight=240mm

%%
%% Три точки для обозначения делимости
%%
\catcode`\@=11
\def\kratno{\mathrel{\smash{\lower.5ex\hbox{$\,\vdots\,$}}}}
% Это прямые три точки для обозначения делимости
\def\nekratno{\mathrel{\mathpalette\c@ncel\kratno}}
\def\c@ncel#1#2{\m@th\ooalign{$\hfil#1\mkern1mu/\hfil$\crcr$#1#2$}}
% Это перечеркнутые три точки для обозначения неделимости
\def\divis{\smash{\lower.1ex\hbox{\,\hbox{.}\kern-.22em\lower-.6ex\hbox{.}
\kern-.52em\lower-1.2ex\hbox{.}\,}} }
% А это наклонные...

\hyphenation{чис-ло чис-ла чис-лу чис-лом чис-ле чис-лам чис-ла-ми чис-лах
че-ты-рех-уголь-ник че-ты-рех-уголь-ни-ка пря-мо-уголь-ник тре-уголь-ни-ка
тре-уголь-ни-ков вне-впи-сан-ной вне-впи-сан-ных}

\sloppy
\pagestyle{empty}

\def\q#1.{{\bf #1. }}
\def\dotline{\smallskip\hbox to \hsize{\dotfill}\medskip}

\begin{document}

\centerline{\sc САНКТ-ПЕТЕРБУРГСКАЯ}
\centerline{\sc ОЛИМПИАДА ШКОЛЬНИКОВ 2023 года ПО МАТЕМАТИКЕ}
\centerline{\sc II тур. \ 9 класс.}
\medskip
\hrule
\bigskip

\q1. Вещественные числа $a$ и $b$, большие 1, удовлетворяют неравенству 
$$
a+\dfrac{1}{a^2} \geq 5b-\dfrac{3}{b^2}.
$$
Докажите, что $a>5b-\dfrac{4}{b^2}$.

\q2. По кругу расставлено несколько (не менее 5) целых чисел таким образом, что каждое из них делится на сумму двух соседних. Сумма всех чисел положительна. Какое наименьшее значение она может принимать? %(22-36 б, СБ)

\q3. Точка $M$~--- середина стороны $AC$ остроугольного треугольника $ABC$. На меньшей дуге $AC$ описанной окружности этого треугольника выбрана точка $K$ таким образом, что угол $AKM$~--- прямой. Отрезки $BK$ и~$AM$ пересекаются в точке~$X$, а~высота, проведенная из вершины~$A$, пересекает отрезок $BM$ в точке $Y$. Докажите, что прямые $XY$ и $AB$ параллельны.

\q4. Одна сторона квадратного листа бумаги покрашена в красный цвет, другая --- в синий. С обеих сторон лист разбит на $300^2$ одинаковых квадратных клеток. В каждой из этих $2 \cdot 300^2$ клеток написано число от 1 до $k$. Найдите наименьшее $k$, при котором могут одновременно выполняться свойства:

\smallskip

\noindent
(i) на красной стороне любые два числа, стоящие в разных строках, различны;
%на красной стороне никакое число не встречается в двух разных строках;

\noindent
(ii) на синей стороне любые два числа, стоящие в разных столбцах, различны;
%на синей стороне никакое число не встречается в двух разных столбцах;

\noindent
(iii) у каждого из $300^2$ квадратов разбиения число на синей стороне не равно числу на красной.

\dotline

\centerline{Олимпиада 2023 года. II тур. 9 класс. Выводная аудитория.}
\smallskip

\q5. Для натурального числа $n$ и ненулевой цифры $d$
обозначим через $f(n,d)$ наименьшее натуральное $k$,
для которого число $kn$ начинается с цифры $d$.
Какое максимальное значение может принимать величина
$f(n,d)$?

\q6. Барон Мюнхгаузен в шутку сказал, что рассадил $2\,000\,000$ кроликов в $1\,000\,000$ клеток так, что в каждой клетке оказалось не более одного кролика. Бургомистр воспринял эти слова всерьёз, и задает барону вопросы вида
<<верно ли, что такой-то кролик сидит в такой-то клетке>>. Бургомистр считает, что он уличит барона во лжи в одном из
трех случаев: 

 1) если барон ответит <<да>> на вопросы про одного и того же кролика и две разные клетки;

 2) если барон ответит <<да>> на вопросы про одну и ту же клетку и двух разных кроликов;

 3) если барон ответит <<нет>> на вопросы про какого-то кролика и все имеющиеся клетки.

Сможет ли барон отвечать так, что бургомистр не сумеет уличить его во лжи, задав $10^{11}$ вопросов?

\q7. Дан треугольник $ABC$. Точка $X$ симметрична вершине $B$ относительно прямой $AC$, а точка $Y$ симметрична $C$ относительно прямой $AB$. Касательная к описанной окружности треугольника $XAY$, проведенная в~точке $A$, пересекает прямые $XY$ и $BC$ в точках $E$ и $F$ соответственно. Докажите, что $AE=AF$.


\clearpage

\centerline{\sc САНКТ-ПЕТЕРБУРГСКАЯ}
\centerline{\sc ОЛИМПИАДА ШКОЛЬНИКОВ 2023 года ПО МАТЕМАТИКЕ}
\centerline{\sc II тур. \ 10 класс.}
\medskip
\hrule
\medskip

\q1. Существует ли набор из 2023 ненулевых вещественных чисел  (не обязательно различных), в котором
дробная часть каждого из этих чисел равна сумме остальных 2022 чисел набора?
   
(Напомним, что дробная часть $\{x\}$ числа $x$ --- это такое число, что $0\leq $\linebreak$\leq\{x\}<1$ и $x-\{x\}$ --- целое число.)

\q2. На биссектрисе угла $B$ треугольника $ABC$, в котором  $\angle B=120^\circ$, отмечена такая точка $D$, что  $\angle ADB=2\angle ACB$. На
отрезке $AB$ отмечена такая точка $E$, что  $AE=AD$. Докажите, что $EC=ED$.

\q3. Найдите все натуральные $x$, $y$ и простые $p$ такие, что $x^5+y^4=pxy$.

\q4. Дано натуральное число $n$. Одна сторона квадратного листа бумаги покрашена в красный цвет, а другая --- в  синий. С обеих сторон лист разбит на~$n^2$ одинаковых квадратных клеток. В каждой из этих $2 n^2$ клеток написано число от 1 до $k$. Найдите наименьшее $k$, при котором могут одновременно выполняться следующие свойства:

\noindent
(i) на красной стороне любые два числа, стоящие в разных строках, различны;

\noindent
(ii) на синей стороне любые два числа, стоящие в разных столбцах, различны;

\noindent
(iii) у каждого из $n^2$ квадратов разбиения число на синей стороне не равно числу на красной.

\dotline

\centerline{Олимпиада 2023 года. II тур. 10 класс. Выводная аудитория.}
\smallskip

\q5. Даны вещественные числа $x_0$, $x_1$, $x_2$, \dots, $x_{n-1}$, $x_n=x_0$
и~функция $f\colon \mathbb{R}\to \mathbb{R}$. При каждом $i=0,1,\dots,n-1$ на отрезке между
точками $x_i$ и $x_{i+1}$ выбирается произвольная точка $y_i$. После этого
вычисляется сумма
$$
(x_1-x_0)f(y_0)+(x_2-x_1)f(y_1)+\ldots+(x_n-x_{n-1})f(y_{n-1}).
$$ 
Докажите, что среди таких сумм есть как неположительные, так и неотрицательные.

\q6. В собрании клуба <<Диоген>> принимают участие несколько джентльменов, некоторые из которых дружат между собой
(дружба взаимна). Назовём участника \emph{нелюдимым}, если у него среди присутствующих на собрании ровно один друг. По
правилам клуба единственный друг любого нелюдимого участника  может покинуть собрание (джентльмены покидают собрание по
одному). Цель собрания --- добиться того, чтобы среди участников не осталось друзей. Докажите, что если цель вообще может быть достигнута, то количество оставшихся на собрании участников не зависит от того, кто и в каком порядке уходил.

\q7. Если $\ell_1$ и $\ell_2$ --- две не параллельные прямые на плоскости, а~$d_1,d_2$ --- положительные
числа, то множество таких точек $X$, что  при $i=1,2$ расстояние от $X$ до $\ell_i$ есть целое кратное
$d_i$, назовём \emph{решёткой}.

Пусть $A$ --- конечное множество точек на плоскости, не лежащих на одной прямой.
Треугольник с вершинами в $A$ назовём \emph{пустым}, если внутри и на сторонах
этого треугольника нет других точек множества $A$, кроме вершин.
Оказалось, что все пустые треугольники с вершинами в $A$ имеют одинаковую площадь. 
Докажите, что существуют решётка $L$ и выпуклый многоугольник $F$ такие, что $A=L\cap F$.


\clearpage

\centerline{\sc САНКТ-ПЕТЕРБУРГСКАЯ}
\centerline{\sc ОЛИМПИАДА ШКОЛЬНИКОВ 2023 года ПО МАТЕМАТИКЕ}
\centerline{\sc II тур. \ 11 класс.}
\medskip
\hrule
\bigskip

\q1. Пусть $f(x)$ --- квадратный трёхчлен, $g(x)$ --- многочлен степени 3. Может ли многочлен $f(g(x))$ иметь шесть различных корней, являющихся степенями~2? 
 
\q2. В треугольнике $ABC$ проведена медиана $BM$. На продолжении стороны $AC$ за точку $C$ отмечена такая точка $D$, что $BD=2CD$. Описанная окружность треугольника $BMC$ пересекает отрезок $BD$ в точке~$N$. Докажите, что $AC+BM>2MN$.
 
\q3. Несократимые дроби $a\over b$ и $c\over d$ записали в виде чисто
периодических десятичных дробей. Оказалось, что любая конечная последовательность
подряд стоящих цифр, встречающаяся в первой десятичной дроби после запятой,
встречается и во второй (тоже подряд и тоже после запятой). Докажите, что $b=d$.

\q4. За какое наименьшее число операций можно перекрасить в чёрный цвет белую доску $100\times 100$, если за одну операцию можно перекрасить в противоположный цвет любые 99 клеток любой горизонтали или любой вертикали?

\dotline

\centerline{Олимпиада 2023 года. II тур. 11 класс. Выводная аудитория.}
\smallskip

\q5. Дано натуральное число $a>1$. Определим функцию
$$
f(n)=n+\bigl[a\{n\sqrt 2\,\}\bigr].
$$
Докажите, что существует такое натуральное $n$, что $f(f(n))=f(n)$, но $f(n)\ne n$. 

Напомним, что $[x]$ --- {\em целая часть} $x$, то есть  наибольшее целое число,  не превосходящее $x$, а $\{x\} = x-[x]$ --- {\em дробная  часть} $x$.

\q6. Дан треугольник $ABC$. Точка $X$ симметрична вершине $B$ относительно прямой $AC$, а точка $Y$ симметрична $C$ относительно $AB$. Касательная к описанной окружности треугольника $XAY$, проведенная в~точке $A$, пересекает прямые $XY$ и $BC$ в точках $E$ и $F$ соответственно. Докажите, что $AE=AF$.

\q7. В связном графе $G$ даны два непересекающихся множества вершин $X$ и $Y$, между этими множествами нет ни одного ребра. Пусть в~графе $G-X$ ровно $m$ компонент связности, а в графе $G-Y$ ровно $n$ компонент связности. Какое наименьшее количество компонент связности может быть в графе  $G-(X\cup Y)$?

Напомним, что для множества вершин $W$ через $G-W$ обозначается граф, полученный из $G$ в результате удаления всех вершин множества $W$ и всех выходящих из них ребер.


\end{document}
